\documentclass[12pt]{article}
\usepackage{float, amsmath, amssymb, amsthm, algorithm, algorithmic, graphicx, caption, subcaption, mathrsfs, color, cancel, verbatim, cite, authblk, mathtools}
\usepackage{enumitem}

\def\upint{\mathchoice%
    {\mkern13mu\overline{\vphantom{\intop}\mkern7mu}\mkern-20mu}%
    {\mkern7mu\overline{\vphantom{\intop}\mkern7mu}\mkern-14mu}%
    {\mkern7mu\overline{\vphantom{\intop}\mkern7mu}\mkern-14mu}%
    {\mkern7mu\overline{\vphantom{\intop}\mkern7mu}\mkern-14mu}%
  \int}
\def\lowint{\mkern3mu\underline{\vphantom{\intop}\mkern7mu}\mkern-10mu\int}

\let\oldemptyset\emptyset
\let\emptyset\varnothing

\setlength{\oddsidemargin}{0in}
\setlength{\evensidemargin}{\oddsidemargin}
\setlength{\textwidth}{6.5in}
\setlength{\topmargin}{-.25in}
\setlength{\headheight}{0in}
\setlength{\headsep}{0in}
\setlength{\topskip}{0in}
\setlength{\textheight}{9.5in}
\font\bigbf = cmbx10 scaled \magstep1
\font\medbf = cmbx10 scaled \magstephalf
\font\medrm = cmr10 scaled \magstephalf
\font\bigrm = cmr10 scaled \magstep1

\usepackage[english]{babel}
\usepackage[utf8]{inputenc}
\usepackage[colorinlistoftodos]{todonotes}

\title{Harmonic Analysis}
\begin{document}
\noindent \textbf{An Introduction to Harmonic Analysis}
\noindent \textbf{Lecture 2: Ricardo S\'aenz} \\
\noindent http://fejer.ucol.mx/ricardo/ \\

\noindent \textbf{Motivation}
\noindent One of the motivations for Harmonic Analysis is diffusion, which are related to problems involving Partial Differential Equations. One common problem is heat diffusion. \\

\noindent Suppose that we are studying the heat diffusion of a square region $S$ (shown in Figure \ref{square}). Suppose that heat is given by the function $u(x,y,t)$, where $(x,y)$ is a coordinate on the square and $t$ is a given time.  

\begin{figure}[H]
\centering
\includegraphics[width=0.2\textwidth]{2Square.pdf}
\caption{}
\label{square}
\end{figure}

\noindent We can describe heat using the following equation, where $\sigma$ is the specific heat,
$$H(t) = \sigma \iint_S u(x,y,t)\,dx\,dy$$
\noindent and we can obtain the flow of heat by taking the partial derivative of $H(t)$ with respect to $t$. This gives us the following flow equation:
\begin{align*}
\frac{dH}{dt} &= \sigma \iint_S \frac{\partial u}{\partial t} (x,y,t) \, dx \, dy \\
&\approx h^2 \cdot \frac{\partial u}{\partial t} (x,y,t)
\end{align*}

\noindent By Newton's Law on the right side: 
$$h\cdot k \Big( \frac{\partial u}{\partial x} (x_0 + \frac{h}{2}, y_0, t) - \frac{\partial u}{\partial x} ( x_0 - \frac{h}{2}, y_0, t) + \frac{\partial u}{\partial y} (x_0, y_0 + \frac{h}{2}, t) - \frac{\partial u}{\partial y} (x_0, y_0-\frac{h}{2},y)\Big) $$
\end{document}